\documentclass[11pt, a4paper]{report}

% ========== Common Packages ==========
\usepackage{fontspec}
\usepackage{kotex}
\usepackage{amsmath, amssymb, amsthm}
\usepackage{mathtools}
\usepackage{bm}
\usepackage{graphicx}
\usepackage{booktabs}
\usepackage{array}
\usepackage{tabularx}
\usepackage{longtable}
\usepackage{hyperref}
\usepackage{cleveref}
\usepackage{geometry}
\usepackage{fancyhdr}
\usepackage{titlesec}
\usepackage{enumitem}
\usepackage{xcolor}
\usepackage{tcolorbox}
\usepackage{algorithm}
\usepackage{algorithmic}
\usepackage{tikz}
\usetikzlibrary{arrows.meta, positioning, calc, decorations.pathreplacing, shapes.geometric, fit, patterns, angles, quotes, trees}
\usepackage{pgfplots}
\pgfplotsset{compat=1.18}
\usepackage{caption}
\usepackage{subcaption}
\usepackage{float}

\geometry{margin=2.5cm, top=3cm, bottom=3cm}

\usepackage{multirow}
% ========== Colors ==========
% Common
\definecolor{defblue}{RGB}{30, 80, 150}
\definecolor{thmgreen}{RGB}{20, 110, 60}
\definecolor{noteorange}{RGB}{200, 100, 0}
\definecolor{lightblue}{RGB}{230, 240, 255}
\definecolor{lightgreen}{RGB}{230, 255, 235}
\definecolor{lightorange}{RGB}{255, 245, 230}
\definecolor{lightgray}{RGB}{245, 245, 245}
% Extended
\definecolor{lightred}{RGB}{255, 235, 235}
\definecolor{darkred}{RGB}{180, 30, 30}
\definecolor{biogreen}{RGB}{10, 130, 80}
\definecolor{cvpurple}{RGB}{100, 40, 150}
\definecolor{injuryred}{RGB}{200, 40, 40}
\definecolor{codegray}{RGB}{240, 240, 245}
\definecolor{pipeblue}{RGB}{25, 100, 180}
\definecolor{constraintred}{RGB}{200, 50, 50}
\definecolor{termblack}{RGB}{30, 30, 30}
\definecolor{goldcolor}{RGB}{180, 140, 20}

% ========== Theorem-like environments ==========
\theoremstyle{definition}
\newtheorem{definition}{Definition}[chapter]
\newtheorem{example}{Example}[chapter]

\theoremstyle{plain}
\newtheorem{theorem}{Theorem}[chapter]
\newtheorem{proposition}{Proposition}[chapter]

\theoremstyle{remark}
\newtheorem{remark}{Remark}[chapter]

% ========== tcolorbox environments ==========
% --- Common (all parts) ---
\newtcolorbox{keyidea}[1][]{
  colback=lightblue,
  colframe=defblue,
  fonttitle=\bfseries,
  title={Key Idea},
  #1
}

\newtcolorbox{importantnote}[1][]{
  colback=lightorange,
  colframe=noteorange,
  fonttitle=\bfseries,
  title={Important},
  #1
}

\newtcolorbox{summary}[1][]{
  colback=lightgreen,
  colframe=thmgreen,
  fonttitle=\bfseries,
  title={Summary},
  #1
}

% --- Warning/Error ---
\newtcolorbox{warningbox}[1][]{
  colback=lightred,
  colframe=darkred,
  fonttitle=\bfseries,
  title={Warning},
  #1
}

% --- Domain: Biomechanics & Clinical ---
\newtcolorbox{clinicalnote}[1][]{
  colback=lightred,
  colframe=darkred,
  fonttitle=\bfseries,
  title={Clinical Note},
  #1
}

\newtcolorbox{biomechbox}[1][]{
  colback=lightgreen,
  colframe=biogreen,
  fonttitle=\bfseries,
  title={Biomechanics},
  #1
}

\newtcolorbox{injurybox}[1][]{
  colback={injuryred!8},
  colframe=injuryred,
  fonttitle=\bfseries,
  title={Injury Risk},
  #1
}

% --- Domain: Computer Vision ---
\newtcolorbox{cvbox}[1][]{
  colback={cvpurple!5},
  colframe=cvpurple,
  fonttitle=\bfseries,
  title={Computer Vision},
  #1
}

% --- Pipeline & Setup ---
\newtcolorbox{setupbox}[1][]{
  colback=lightgray,
  colframe=gray!60,
  fonttitle=\bfseries,
  title={Setup},
  #1
}

\newtcolorbox{pipelinebox}[1][]{
  colback=pipeblue!5,
  colframe=pipeblue,
  fonttitle=\bfseries,
  title={Pipeline Step},
  #1
}

\newtcolorbox{codebox}[1][]{
  colback=codegray,
  colframe=gray!70,
  fonttitle=\bfseries\ttfamily,
  title={Code},
  #1
}

\newtcolorbox{termbox}[1][]{
  colback=termblack!5,
  colframe=termblack!60,
  fonttitle=\bfseries\ttfamily,
  title={Terminal},
  #1
}

% --- Research ---
\newtcolorbox{hypothesis}[1][]{
  colback=goldcolor!8,
  colframe=goldcolor,
  fonttitle=\bfseries,
  title={Hypothesis},
  #1
}

\newtcolorbox{resultbox}[1][]{
  colback=thmgreen!8,
  colframe=thmgreen,
  fonttitle=\bfseries,
  title={Expected Result},
  #1
}

% --- Constraints & Solutions ---
\newtcolorbox{constraintbox}[1][]{
  colback=constraintred!8,
  colframe=constraintred,
  fonttitle=\bfseries,
  title={Constraint},
  #1
}

\newtcolorbox{solutionbox}[1][]{
  colback=thmgreen!8,
  colframe=thmgreen,
  fonttitle=\bfseries,
  title={Solution},
  #1
}

\newtcolorbox{databox}[1][]{
  colback=pipeblue!5,
  colframe=pipeblue,
  fonttitle=\bfseries,
  title={Data Spec},
  #1
}

% --- Progress/Status ---
\newtcolorbox{successbox}[1][]{
  colback=lightgreen,
  colframe=thmgreen,
  fonttitle=\bfseries,
  title={Success},
  #1
}

\newtcolorbox{nextbox}[1][]{
  colback=lightorange,
  colframe=noteorange,
  fonttitle=\bfseries,
  title={Next Step},
  #1
}

% --- Fitting guide ---
\newtcolorbox{failbox}[1][]{
  colback=lightred,
  colframe=darkred,
  fonttitle=\bfseries,
  title={Failure Pattern},
  #1
}

\newtcolorbox{fixbox}[1][]{
  colback=lightgreen,
  colframe=thmgreen,
  fonttitle=\bfseries,
  title={Fix},
  #1
}

\newtcolorbox{lessonbox}[1][]{
  colback=lightorange,
  colframe=noteorange,
  fonttitle=\bfseries,
  title={Lesson Learned},
  #1
}

% --- Specification ---
\newtcolorbox{specbox}[1][]{
  colback=cvpurple!5,
  colframe=cvpurple,
  fonttitle=\bfseries,
  title={Specification},
  #1
}

\newtcolorbox{checklistbox}[1][]{
  colback=lightgreen,
  colframe=biogreen,
  fonttitle=\bfseries,
  title={Checklist},
  #1
}

% ========== Custom math commands ==========
% Vectors (lowercase)
\newcommand{\vx}{\mathbf{x}}
\newcommand{\vd}{\mathbf{d}}
\newcommand{\vo}{\mathbf{o}}
\newcommand{\vc}{\mathbf{c}}
\newcommand{\vr}{\mathbf{r}}
\newcommand{\vp}{\mathbf{p}}
\newcommand{\vv}{\mathbf{v}}
\newcommand{\vn}{\mathbf{n}}
\newcommand{\vu}{\mathbf{u}}
\newcommand{\vt}{\mathbf{t}}
\newcommand{\ve}{\mathbf{e}}
\newcommand{\vq}{\mathbf{q}}
% Vectors (uppercase)
\newcommand{\vF}{\mathbf{F}}
\newcommand{\vM}{\mathbf{M}}
\newcommand{\va}{\mathbf{a}}
\newcommand{\vg}{\mathbf{g}}
\newcommand{\vX}{\mathbf{X}}
% Bold Greek
\newcommand{\vmu}{\bm{\mu}}
\newcommand{\vbeta}{\bm{\beta}}
\newcommand{\vtheta}{\bm{\theta}}
\newcommand{\vpsi}{\bm{\psi}}
% Matrices
\newcommand{\mSigma}{\bm{\Sigma}}
\newcommand{\mR}{\mathbf{R}}
\newcommand{\mS}{\mathbf{S}}
\newcommand{\mW}{\mathbf{W}}
\newcommand{\mJ}{\mathbf{J}}
\newcommand{\mG}{\mathbf{G}}
\newcommand{\mT}{\mathbf{T}}
\newcommand{\mI}{\mathbf{I}}
\newcommand{\mB}{\mathbf{B}}
\newcommand{\mK}{\mathbf{K}}
\newcommand{\mP}{\mathbf{P}}
\newcommand{\mF}{\mathbf{F}}
\newcommand{\mE}{\mathbf{E}}
\newcommand{\mA}{\mathbf{A}}
\newcommand{\mH}{\mathbf{H}}
% Sets and groups
\newcommand{\RR}{\mathbb{R}}
\newcommand{\SO}{\mathrm{SO}}
% Units
\newcommand{\degps}{\,\text{deg/s}}
\newcommand{\Nm}{\ensuremath{\,\text{N}\cdot\text{m}}}
\newcommand{\BW}{\,\text{BW}}

% ========== Listings ==========
\usepackage{listings}

\lstset{
  basicstyle=\ttfamily\small,
  keywordstyle=\color{defblue}\bfseries,
  commentstyle=\color{thmgreen}\itshape,
  stringstyle=\color{noteorange},
  backgroundcolor=\color{codegray},
  frame=single,
  rulecolor=\color{gray!40},
  breaklines=true,
  showstringspaces=false,
  numbers=left,
  numberstyle=\tiny\color{gray},
  tabsize=4,
  captionpos=b
}

% ========== Header/Footer ==========
\pagestyle{fancy}
\fancyhf{}
\fancyhead[L]{\leftmark}
\fancyhead[R]{\thepage}
\renewcommand{\headrulewidth}{0.4pt}

% ========== Title formatting ==========
\titleformat{\chapter}[display]
  {\normalfont\huge\bfseries\color{defblue}}
  {\chaptertitlename\ \thechapter}{20pt}{\Huge}

\titleformat{\section}
  {\normalfont\Large\bfseries\color{defblue!80!black}}
  {\thesection}{1em}{}

\titleformat{\subsection}
  {\normalfont\large\bfseries\color{defblue!60!black}}
  {\thesubsection}{1em}{}


\begin{document}

% ==================== TITLE PAGE ====================
\begin{titlepage}
\centering
\vspace*{1.5cm}
{\Large\color{gray} Part 4 --- Fitting Guide}
\vspace{0.5cm}

{\Huge\bfseries\color{defblue} OpenPose 3D 키포인트에서\\[0.3cm]
SMPL 인체 모델 피팅까지}

\vspace{1cm}

{\LARGE\color{defblue!70!black} 3D Keypoints $\rightarrow$ SMPL Parameters:\\
수학적 정식화, 단계별 파이프라인, 실전 구현}

\vspace{2cm}

\begin{tikzpicture}[
  block/.style={rectangle, rounded corners=4pt, draw=defblue, fill=lightblue,
    minimum width=2.2cm, minimum height=0.7cm, font=\small\bfseries, text=defblue!80!black, align=center},
  arrow/.style={-Stealth, thick, defblue!60}
]
\node[block] (op2d) at (0,0) {OpenPose\\2D $\times$ 7};
\node[block] (tri) at (3.5,0) {DLT\\삼각측량};
\node[block, fill=cvpurple!15, draw=cvpurple] (op3d) at (7,0) {OpenPose\\3D (25J)};
\node[block, fill=noteorange!15, draw=noteorange] (smpl) at (10.5,0) {SMPL\\피팅};
\node[block, fill=thmgreen!15, draw=thmgreen] (mesh) at (14,0) {SMPL\\Mesh};

\draw[arrow] (op2d) -- (tri);
\draw[arrow] (tri) -- (op3d);
\draw[arrow, noteorange] (op3d) -- (smpl) node[midway, above, font=\tiny\color{noteorange}] {이 문서};
\draw[arrow, thmgreen] (smpl) -- (mesh);
\end{tikzpicture}

\vspace{2cm}
{\large 2026}

\vspace{\fill}
{\small\color{gray} 이 문서는 Part 1 (SMPL 모델 이론), Part 2 (동작 분석), Part 3 (투구 파이프라인)의 후속으로,\\
OpenPose 3D 키포인트에서 SMPL 파라미터를 최적화하는 전체 과정을 상세히 다룹니다.}
\end{titlepage}

\tableofcontents
\newpage


% ====================================================================
% CHAPTER 1
% ====================================================================
\chapter{OpenPose 3D 키포인트의 이해}
\label{ch:openpose3d}

\section{OpenPose Body25 관절 정의}

OpenPose Body25 모델은 25개의 관절(keypoint)을 정의한다.

\begin{table}[H]
\centering
\caption{OpenPose Body25 관절 인덱스}
\renewcommand{\arraystretch}{1.2}
\small
\begin{tabular}{cl|cl|cl}
\toprule
\textbf{\#} & \textbf{이름} & \textbf{\#} & \textbf{이름} & \textbf{\#} & \textbf{이름} \\
\midrule
0 & Nose & 9 & RHip & 18 & LEar \\
1 & Neck & 10 & RKne & 19 & LBToe \\
2 & RSho & 11 & RAnk & 20 & LSToe \\
3 & RElb & 12 & LHip & 21 & LHeel \\
4 & RWri & 13 & LKne & 22 & RBToe \\
5 & LSho & 14 & LAnk & 23 & RSToe \\
6 & LElb & 15 & REye & 24 & RHeel \\
7 & LWri & 16 & LEye & & \\
8 & MidHip & 17 & REar & & \\
\bottomrule
\end{tabular}
\end{table}


\section{2D에서 3D로: 다시점 삼각측량}

$C$개 카메라에서 관절 $j$의 2D 좌표 $\mathbf{u}_{cj} = (x_{cj}, y_{cj})$가 주어질 때,
투영 행렬 $\mK_c[\mR_c | \vt_c]$를 이용한 DLT 삼각측량:

\begin{equation}
\mathbf{X}_j^* = \arg\min_{\mathbf{X} \in \RR^3}
  \sum_{c=1}^{C} w_{cj} \left\| \mathbf{u}_{cj} - \pi_c(\mathbf{X}) \right\|^2
\end{equation}

신뢰도 $w_{cj}$를 가중치로 사용하여, 검출이 확실한 뷰에 더 큰 비중을 둔다.
최소 3개 뷰 이상에서 검출된 관절만 삼각측량한다.


\section{OpenPose $\leftrightarrow$ SMPL 관절 매핑}

\begin{keyidea}
OpenPose Body25 (25개)와 SMPL (24개)의 관절 정의가 다르므로,
\textbf{대응 관절 매핑}이 피팅의 전제조건이다.
얼굴 관절(Eyes, Ears)과 발가락 관절(BToe, SToe, Heel)은 SMPL에 직접 대응이 없다.
\end{keyidea}

\begin{table}[H]
\centering
\caption{OpenPose Body25 $\rightarrow$ SMPL 24 관절 매핑 (14개 대응)}
\label{tab:joint_mapping}
\renewcommand{\arraystretch}{1.3}
\begin{tabular}{clcl}
\toprule
\textbf{OP25 \#} & \textbf{OP 이름} & \textbf{SMPL \#} & \textbf{SMPL 이름} \\
\midrule
1 & Neck & 12 & Neck \\
8 & MidHip & 0 & Pelvis \\
2 & RSho & 17 & R\_Shoulder \\
3 & RElb & 19 & R\_Elbow \\
4 & RWri & 21 & R\_Wrist \\
5 & LSho & 16 & L\_Shoulder \\
6 & LElb & 18 & L\_Elbow \\
7 & LWri & 20 & L\_Wrist \\
9 & RHip & 2 & R\_Hip \\
10 & RKne & 5 & R\_Knee \\
11 & RAnk & 8 & R\_Ankle \\
12 & LHip & 1 & L\_Hip \\
13 & LKne & 4 & L\_Knee \\
14 & LAnk & 7 & L\_Ankle \\
\midrule
\multicolumn{4}{l}{\textcolor{gray}{매핑 불가: 0(Nose), 15--18(Eyes/Ears), 19--24(Toes/Heel)}} \\
\bottomrule
\end{tabular}
\end{table}


\section{좌표계 정합}

\begin{importantnote}
OpenPose 3D와 SMPL의 좌표계가 다를 수 있다.
\begin{itemize}[nosep]
  \item \textbf{OpenPose 3D}: 카메라 보정의 월드 좌표계 (DA3 PnP 등)
  \item \textbf{SMPL}: 정규 공간에서 정의, 글로벌 회전 $\mR$과 이동 $\vt$로 변환
\end{itemize}
피팅 과정에서 $\mR$과 $\vt$를 최적화하여 두 좌표계를 자동으로 정합한다.
별도의 좌표 변환은 불필요하다.
\end{importantnote}


% ====================================================================
% CHAPTER 2
% ====================================================================
\chapter{SMPL 피팅의 수학적 정식화}
\label{ch:formulation}

\section{문제 정의}

프레임 $t$에서 OpenPose 3D 키포인트 $\{\vp_j^{\text{OP}}\}_{j \in \mathcal{M}}$이 주어질 때,
SMPL 파라미터를 최적화한다:

\begin{equation}
\boxed{
\vtheta^*, \vbeta^*, \mR^*, \vt^* = \arg\min_{\vtheta, \vbeta, \mR, \vt}
E_{\text{total}}(\vtheta, \vbeta, \mR, \vt \mid \{\vp_j^{\text{OP}}\})
}
\label{eq:total_objective}
\end{equation}

여기서:
\begin{itemize}[nosep]
  \item $\vtheta \in \RR^{72}$: body pose (24 관절 $\times$ 3 axis-angle)
  \item $\vbeta \in \RR^{10}$: body shape (PCA 계수)
  \item $\mR \in SO(3)$: 글로벌 회전 (root orientation)
  \item $\vt \in \RR^3$: 글로벌 이동 (root translation)
\end{itemize}


\section{에너지 함수}

\begin{equation}
E_{\text{total}} = \lambda_{J} E_{J3D} + \lambda_\theta E_\theta + \lambda_\beta E_\beta
  + \lambda_{\text{smooth}} E_{\text{smooth}} + \lambda_{\text{limit}} E_{\text{limit}}
\label{eq:energy}
\end{equation}

\subsection{3D 관절 위치 오차 ($E_{J3D}$)}

핵심 데이터 항. SMPL의 관절 위치와 OpenPose 3D 키포인트 사이의 오차를 최소화한다:

\begin{equation}
E_{J3D} = \sum_{j \in \mathcal{M}} w_j \rho\!\left(
  \left\| J_j^{\text{SMPL}}(\vtheta, \vbeta) - \vp_j^{\text{OP}} \right\|^2
\right)
\end{equation}

여기서:
\begin{itemize}[nosep]
  \item $J_j^{\text{SMPL}}(\vtheta, \vbeta)$: SMPL forward kinematics로 계산된 관절 $j$의 3D 위치
  \item $\vp_j^{\text{OP}}$: OpenPose 3D 삼각측량 결과
  \item $w_j$: 관절별 가중치 (삼각측량 신뢰도 기반)
  \item $\rho(\cdot)$: 로버스트 비용 함수
  \item $\mathcal{M}$: 매핑 가능한 관절 집합 (\cref{tab:joint_mapping}의 14개)
\end{itemize}

\begin{cvbox}[title={로버스트 비용 함수}]
이상치에 강건한 Geman-McClure 비용 함수를 권장한다:
\begin{equation}
\rho_{\text{GM}}(r) = \frac{r}{\sigma^2 + r}
\end{equation}
$\sigma$는 이상치 임계값으로, 일반적으로 $\sigma = 100$\,mm.
큰 오차에 대해 기울기가 포화되어 이상치의 영향을 억제한다.
\end{cvbox}

\begin{keyidea}
\textbf{SMPLify와의 핵심 차이}:
SMPLify는 2D 키포인트를 사용하므로 재투영 오차 $\|\pi(\cdot) - \mathbf{u}\|^2$를 최소화하고,
카메라 파라미터가 필요하다.
\textbf{3D 키포인트를 직접 사용하면 재투영 과정이 불필요하고, 카메라 파라미터에 독립적이다.}
이것이 3D 피팅의 가장 큰 장점이다.
\end{keyidea}


\subsection{자세 사전 ($E_\theta$)}

해부학적으로 불가능한 자세를 방지하는 사전(prior):

\begin{equation}
E_\theta = \|\vtheta\|^2 \quad \text{(간단한 $L_2$ 정규화)}
\end{equation}

또는 학습된 자세 공간을 사용:
\begin{equation}
E_\theta^{\text{VPoser}} = \|z_\theta\|^2, \quad \vtheta = \text{VPoser}(z_\theta)
\end{equation}
VPoser는 VAE로 학습된 인체 자세의 잠재 공간으로,
$z_\theta \in \RR^{32}$에서 $\vtheta \in \RR^{63}$ (21 관절, root 제외)을 생성한다.

\begin{importantnote}
$L_2$ 정규화는 모든 관절을 T-pose로 당기는 경향이 있다.
투구와 같은 극단적 자세에서는 $\lambda_\theta$를 작게 설정해야 한다.
VPoser를 사용하면 이 문제를 완화할 수 있지만, 사전 학습 데이터에 투구 동작이 포함되어 있지 않을 수 있다.
\end{importantnote}


\subsection{체형 정규화 ($E_\beta$)}

\begin{equation}
E_\beta = \|\vbeta\|^2
\end{equation}

체형이 평균(zero)에서 크게 벗어나지 않도록 한다.
같은 피험자의 시퀀스에서는 $\vbeta$를 모든 프레임에서 공유한다.


\subsection{시간적 매끄러움 ($E_{\text{smooth}}$)}

시퀀스 데이터의 경우 프레임 간 급격한 변화를 억제한다:

\begin{equation}
E_{\text{smooth}} = \sum_t \left( \|\vtheta_t - \vtheta_{t-1}\|^2
  + \gamma \|\vtheta_t - 2\vtheta_{t-1} + \vtheta_{t-2}\|^2 \right)
\end{equation}

첫째 항은 속도 정규화, 둘째 항은 가속도 정규화이다.

\begin{biomechbox}[title={투구 동작에서의 시간 정규화}]
투구 가속 단계(약 30\,ms)에서는 어깨 내회전이 7,000$^\circ$/s에 달한다.
이 구간에서 시간 정규화를 과도하게 적용하면 실제 빠른 동작을 억제한다.
Part 5의 GART-Pitch에서 제안한 \textbf{위상별 적응형 $\lambda_{\text{smooth}}$}를 적용할 수 있다.
\end{biomechbox}


\subsection{관절 각도 제한 ($E_{\text{limit}}$)}

해부학적으로 불가능한 관절 각도를 소프트 페널티로 방지:

\begin{equation}
E_{\text{limit}} = \sum_{k \in \mathcal{J}} \left[
  \max(0, \theta_k - \theta_k^{\max})^2 + \max(0, \theta_k^{\min} - \theta_k)^2
\right]
\end{equation}

예: 팔꿈치 과신전 $< 5^\circ$, 무릎 과신전 $< 0^\circ$.


\section{Coarse-to-Fine 최적화 전략}

\begin{table}[H]
\centering
\caption{Coarse-to-Fine 최적화 단계}
\renewcommand{\arraystretch}{1.3}
\begin{tabular}{clcccc}
\toprule
\textbf{Stage} & \textbf{최적화 변수} & \textbf{고정} & \textbf{학습률} & \textbf{Iterations} \\
\midrule
1 & $\mR, \vt$ & $\vtheta, \vbeta$ & $10^{-2}$ & 100 \\
2 & $\vbeta$ & $\mR, \vt, \vtheta$ & $10^{-2}$ & 200 \\
3 & $\vtheta$ & $\mR, \vt, \vbeta$ & $10^{-3}$ & 500 \\
4 & $\vtheta, \vbeta, \mR, \vt$ & --- & $10^{-4}$ & 300 \\
\bottomrule
\end{tabular}
\end{table}


% ====================================================================
% CHAPTER 3
% ====================================================================
\chapter{단계별 피팅 파이프라인}
\label{ch:pipeline}

\begin{figure}[H]
\centering
\begin{tikzpicture}[
  block/.style={rectangle, rounded corners=3pt, draw, minimum width=3.5cm, minimum height=0.8cm,
    align=center, font=\small},
  arrow/.style={-{Stealth[length=2mm]}, thick}
]
\node[block, fill=lightgray] (pre) at (0,0) {Step 0: 전처리\\(필터링, 좌표계)};
\node[block, fill=lightorange] (init) at (0,-1.5) {Step 1: 글로벌 자세\\($\mR, \vt$ 초기화)};
\node[block, fill=lightblue] (beta) at (0,-3) {Step 2: 체형 추정\\($\vbeta$ 최적화)};
\node[block, fill=cvpurple!15] (theta) at (0,-4.5) {Step 3: 자세 최적화\\($\vtheta$ 최적화)};
\node[block, fill=defblue!15] (joint) at (0,-6) {Step 4: 공동 정제\\($\vtheta, \vbeta, \mR, \vt$)};
\node[block, fill=lightgreen] (temp) at (0,-7.5) {Step 5: 시간적 매끄러움\\(시퀀스의 경우)};

\draw[arrow] (pre) -- (init);
\draw[arrow] (init) -- (beta);
\draw[arrow] (beta) -- (theta);
\draw[arrow] (theta) -- (joint);
\draw[arrow] (joint) -- (temp);

% Annotations
\node[font=\scriptsize\color{gray}, right=0.5cm] at (pre.east) {이상치 제거, confidence 필터};
\node[font=\scriptsize\color{gray}, right=0.5cm] at (init.east) {골반 위치 + 몸통 방향};
\node[font=\scriptsize\color{gray}, right=0.5cm] at (beta.east) {사지 길이 비율로 초기 추정};
\node[font=\scriptsize\color{gray}, right=0.5cm] at (theta.east) {관절별 순차 또는 동시};
\node[font=\scriptsize\color{gray}, right=0.5cm] at (joint.east) {Fine-tuning, 학습률 감소};
\node[font=\scriptsize\color{gray}, right=0.5cm] at (temp.east) {가속도 정규화};
\end{tikzpicture}
\caption{SMPL 피팅 5단계 파이프라인}
\end{figure}


\section{Step 0: 전처리}

\begin{enumerate}[leftmargin=*]
  \item \textbf{신뢰도 필터링}: 삼각측량 신뢰도가 낮은 관절 ($w_j < 0.3$) 제거
  \item \textbf{이상치 제거}: 관절 위치가 인체 크기 대비 비현실적인 경우 제거
        (예: 팔 길이 $> 1$\,m)
  \item \textbf{시간적 필터링}: Butterworth 저역통과 필터 (cutoff $\sim$15\,Hz) 적용하여 고주파 노이즈 제거
  \item \textbf{관절 매핑}: OpenPose 25 $\rightarrow$ SMPL 14 대응 관절 추출 (\cref{tab:joint_mapping})
\end{enumerate}


\section{Step 1: 글로벌 자세 초기화}

SMPL의 root를 OpenPose 3D의 골반(MidHip) 위치에 정합한다.

\subsection{Translation 초기화}

\begin{equation}
\vt^{(0)} = \vp_8^{\text{OP}} \quad (\text{MidHip} = \text{Pelvis 근사})
\end{equation}

\subsection{Rotation 초기화}

몸통 방향 벡터로 글로벌 회전을 추정한다:

\begin{align}
\mathbf{up} &= \frac{\vp_1^{\text{OP}} - \vp_8^{\text{OP}}}{\|\vp_1 - \vp_8\|} \quad (\text{Neck} - \text{MidHip}) \\
\mathbf{right} &= \frac{\vp_2^{\text{OP}} - \vp_5^{\text{OP}}}{\|\vp_2 - \vp_5\|} \quad (\text{RSho} - \text{LSho}) \\
\mathbf{front} &= \mathbf{up} \times \mathbf{right}
\end{align}

이로부터 초기 회전 행렬 $\mR^{(0)} = [\mathbf{right} \;|\; \mathbf{up} \;|\; \mathbf{front}]$를 구성한다.


\section{Step 2: 체형 ($\vbeta$) 추정}

사지 길이 비율을 이용한 초기 추정:

\begin{equation}
L_k^{\text{OP}} = \|\vp_{j_1}^{\text{OP}} - \vp_{j_2}^{\text{OP}}\| \quad (k\text{번째 사지})
\end{equation}

SMPL의 $\vbeta$에 따른 사지 길이 $L_k^{\text{SMPL}}(\vbeta)$와의 차이를 최소화:

\begin{equation}
\vbeta^{(0)} = \arg\min_{\vbeta} \sum_k \left( L_k^{\text{SMPL}}(\vbeta) - L_k^{\text{OP}} \right)^2 + \lambda_\beta \|\vbeta\|^2
\end{equation}

시퀀스의 경우, 여러 프레임의 사지 길이 중앙값을 사용하여 안정적으로 추정한다.


\section{Step 3: 자세 ($\vtheta$) 최적화}

\begin{equation}
\vtheta^* = \arg\min_{\vtheta} \sum_{j \in \mathcal{M}} w_j \| J_j^{\text{SMPL}}(\vtheta, \vbeta^{(0)}) + \vt^{(0)} - \vp_j^{\text{OP}} \|^2
  + \lambda_\theta \|\vtheta\|^2
\end{equation}

\begin{importantnote}
SMPL의 forward kinematics는 관절 체인이므로,
\textbf{부모 관절의 오차가 자식 관절로 전파}된다.
근위(proximal) 관절부터 최적화하는 것이 안정적이다:
\begin{enumerate}[nosep]
  \item 골반(root) $\rightarrow$ 척추(spine) $\rightarrow$ 목(neck)
  \item 어깨 $\rightarrow$ 팔꿈치 $\rightarrow$ 손목
  \item 엉덩이 $\rightarrow$ 무릎 $\rightarrow$ 발목
\end{enumerate}
\end{importantnote}


\section{Step 4: 공동 정제}

모든 파라미터를 작은 학습률로 동시 최적화:

\begin{equation}
\vtheta^*, \vbeta^*, \mR^*, \vt^* = \arg\min \;
  \lambda_J E_{J3D} + \lambda_\theta E_\theta + \lambda_\beta E_\beta + \lambda_{\text{limit}} E_{\text{limit}}
\end{equation}

Adam 또는 L-BFGS 옵티마이저를 사용한다.
L-BFGS가 수렴이 빠르지만, 미니배치 처리에는 Adam이 적합하다.


\section{Step 5: 시간적 매끄러움}

시퀀스 전체에 대해 시간 정규화를 추가하여 최종 정제:

\begin{equation}
\min_{\{\vtheta_t\}} \sum_t \left[
  E_{J3D}(\vtheta_t) + \lambda_{\text{smooth}} \left(
    \|\vtheta_t - \vtheta_{t-1}\|^2 + \gamma \|\vtheta_t - 2\vtheta_{t-1} + \vtheta_{t-2}\|^2
  \right)
\right]
\end{equation}

이때 $\vbeta$는 시퀀스 전체에서 공유(단일 값)하고, $\vtheta_t$만 프레임별로 최적화한다.


% ====================================================================
% CHAPTER 4
% ====================================================================
\chapter{실전 구현}
\label{ch:implementation}

\section{PyTorch 구현 의사코드}

\begin{codebox}[title={OpenPose 3D $\rightarrow$ SMPL 피팅 (PyTorch)}]
\begin{lstlisting}[language=Python]
import torch
import smplx

# 1. Load SMPL model
body_model = smplx.create(
    model_path='models/smpl',
    model_type='smpl',
    gender='neutral',
    batch_size=1
).cuda()

# 2. Load OpenPose 3D keypoints
op3d = load_openpose3d('openpose3d_data.json')  # (T, 25, 3)
target = op3d[:, OP_TO_SMPL_IDX]  # (T, 14, 3)

# 3. Initialize parameters
global_orient = torch.zeros(T, 3, requires_grad=True)
body_pose = torch.zeros(T, 69, requires_grad=True)
betas = torch.zeros(1, 10, requires_grad=True)
transl = torch.tensor(op3d[:, 8], requires_grad=True)

# 4. Stage 1: Global orient + translation
opt1 = torch.optim.Adam([global_orient, transl], lr=1e-2)
for i in range(100):
    output = body_model(
        global_orient=global_orient,
        body_pose=body_pose,
        betas=betas.expand(T, -1),
        transl=transl
    )
    joints = output.joints[:, SMPL_TO_OP_IDX]
    loss = ((joints - target) ** 2).sum()
    loss.backward(); opt1.step(); opt1.zero_grad()

# 5. Stage 2: Body shape
opt2 = torch.optim.Adam([betas], lr=1e-2)
for i in range(200):
    ...  # similar loop, optimize betas

# 6. Stage 3: Body pose
opt3 = torch.optim.Adam([body_pose], lr=1e-3)
for i in range(500):
    ...  # + pose prior + joint limits

# 7. Stage 4: Joint refinement
opt4 = torch.optim.Adam(
    [global_orient, body_pose, betas, transl],
    lr=1e-4
)
for i in range(300):
    ...  # all parameters jointly
\end{lstlisting}
\end{codebox}

\section{기존 도구 비교}

\begin{table}[H]
\centering
\caption{3D 키포인트 $\rightarrow$ SMPL 피팅 도구 비교}
\renewcommand{\arraystretch}{1.3}
\begin{tabularx}{\textwidth}{lXcc}
\toprule
\textbf{도구} & \textbf{설명} & \textbf{3D 입력} & \textbf{GPU} \\
\midrule
SMPLify-X & 공식 피팅 도구, 2D/3D 지원 & 가능 & 권장 \\
EasyMocap & 다시점 2D $\rightarrow$ SMPL, 3D 모드 있음 & 부분 & 권장 \\
직접 구현 & smplx + PyTorch 최적화 & \textbf{완전} & 필수 \\
\bottomrule
\end{tabularx}
\end{table}

\begin{keyidea}
이미 정확한 3D 키포인트가 있는 경우, \textbf{직접 구현이 가장 단순하고 확실}하다.
EasyMocap은 2D에서 3D를 재구성하는 과정이 추가되어 오히려 오차가 누적될 수 있다
(Part 0 보고서에서 확인한 좌표계 불일치 문제).
\end{keyidea}


\section{흔한 실패와 디버깅}

\begin{table}[H]
\centering
\caption{SMPL 피팅 실패 유형과 해결}
\renewcommand{\arraystretch}{1.4}
\begin{tabularx}{\textwidth}{lXX}
\toprule
\textbf{증상} & \textbf{원인} & \textbf{해결} \\
\midrule
자세가 뒤집어짐 & 글로벌 회전 초기화 실패 & Step 1의 몸통 방향 초기화 정확히 수행 \\
팔다리 교차 & $\vtheta$ 좌/우 혼동 & 관절 매핑의 L/R 확인, 좌표계 방향 확인 \\
$\vbeta$ 발산 & $\lambda_\beta$ 너무 작음 & $\lambda_\beta$ 증가, 시퀀스 공유 $\vbeta$ 사용 \\
T-pose로 수렴 & $\lambda_\theta$ 너무 큼 & $\lambda_\theta$ 감소, VPoser 사용 \\
떨림(jitter) & 시간 정규화 부재 & Step 5 적용, $\lambda_{\text{smooth}}$ 조절 \\
극단적 자세 실패 & 자세 사전이 투구 미지원 & $\lambda_\theta$ 극소화, 관절 제한만 유지 \\
\bottomrule
\end{tabularx}
\end{table}


% ====================================================================
% CHAPTER 5
% ====================================================================
\chapter{다시점 시나리오에서의 통합}
\label{ch:multiview}

\section{현재 파이프라인에서의 위치}

\begin{figure}[H]
\centering
\begin{tikzpicture}[
  block/.style={rectangle, rounded corners=3pt, draw, minimum width=3cm, minimum height=0.7cm,
    align=center, font=\small},
  arrow/.style={-{Stealth[length=2mm]}, thick},
  done/.style={block, fill=thmgreen!15, draw=thmgreen},
  curr/.style={block, fill=noteorange!15, draw=noteorange, very thick},
  next/.style={block, fill=lightgray, draw=gray}
]
\node[done] (cam) at (0,0) {카메라 보정\\(DA3 PnP)};
\node[done] (seg) at (4,0) {SAM3\\세그멘테이션};
\node[done] (op2d) at (8,0) {OpenPose 2D\\$\times$ 7 카메라};
\node[done] (op3d) at (4,-1.5) {DLT 삼각측량\\OpenPose 3D};
\node[curr] (smpl) at (4,-3) {SMPL 피팅\\$\vtheta, \vbeta, \mR, \vt$};
\node[next] (gart) at (4,-4.5) {GART 학습\\3DGS 인체 재구성};
\node[next] (bio) at (4,-6) {생체역학 분석\\MER, 외반 토크};

\draw[arrow, thmgreen] (cam) -- (seg);
\draw[arrow, thmgreen] (seg) -- (op2d);
\draw[arrow, thmgreen] (op2d) -- (op3d);
\draw[arrow, noteorange, very thick] (op3d) -- (smpl);
\draw[arrow, gray] (smpl) -- (gart);
\draw[arrow, gray] (gart) -- (bio);
\end{tikzpicture}
\caption{GART-Pitch 파이프라인에서 SMPL 피팅의 위치}
\end{figure}


\section{삼각측량 후 피팅 vs 직접 다시점 피팅}

\begin{table}[H]
\centering
\caption{두 가지 다시점 SMPL 피팅 접근법 비교}
\renewcommand{\arraystretch}{1.3}
\begin{tabularx}{\textwidth}{lXX}
\toprule
& \textbf{삼각측량 $\rightarrow$ 3D 피팅} & \textbf{직접 다시점 2D 피팅} \\
\midrule
입력 & 3D 키포인트 (삼각측량 결과) & 2D 키포인트 $\times$ $C$개 뷰 \\
장점 & 카메라 파라미터 불필요, 단순 & 가려짐에 강건, 깊이 모호성 적음 \\
단점 & 삼각측량 오차 누적 & 카메라 파라미터 필요, 복잡 \\
도구 & 직접 구현 (이 문서) & EasyMocap, SMPLify-X \\
\textbf{현재 선택} & \textcolor{thmgreen}{\textbf{채택}} & 좌표계 불일치 문제 \\
\bottomrule
\end{tabularx}
\end{table}


\section{SMPL 피팅 결과의 활용}

피팅된 SMPL 파라미터는 두 가지 경로로 활용된다:

\begin{enumerate}[leftmargin=*]
  \item \textbf{GART 학습 입력}: 프레임별 $\vtheta_t$로 캐노니컬 Gaussian을 자세 변환 (LBS)
  \item \textbf{생체역학 분석}: $\vtheta_t$에서 ISB 관절 각도 추출 $\rightarrow$ MER, 외반 토크 등 산출
\end{enumerate}

\begin{summary}[title={OpenPose 3D $\rightarrow$ SMPL 피팅 핵심 요약}]
\begin{enumerate}[nosep]
  \item \textbf{매핑}: OpenPose 25관절 중 14개가 SMPL 24관절에 대응
  \item \textbf{에너지}: $E_{J3D}$ (3D 관절 오차) + $E_\theta$ (자세 사전) + $E_\beta$ (체형) + $E_{\text{smooth}}$ (시간)
  \item \textbf{최적화}: Coarse-to-fine 4단계 ($\mR,\vt \rightarrow \vbeta \rightarrow \vtheta \rightarrow$ 공동)
  \item \textbf{장점}: 카메라 파라미터 불필요, 좌표계 독립, 직접 3D 감독
  \item \textbf{구현}: smplx + PyTorch로 GPU 서버에서 프레임당 $\sim$1초
\end{enumerate}
\end{summary}


\end{document}
